
Graphviz es una herramienta para hacer diagramas

\url{https://graphviz.org/}

Éstos los podemos integrar en nuestros documentos de manera sencilla:

\digraph{DiagramaSencillo}{
  rankdir=LR;
  a -> b -> c;
  c -> a;
}

Uno más elaborado obtenido de la documentación oficial (\ref{fig:graficoDot} \nameref{fig:graficoDot}).

Puedes probar tus diagramas en \url{https://dreampuf.github.io/GraphvizOnline}

\begin{figure}
    \centering
    \digraph{MasElaborado}{
      subgraph cluster_0 {
        style=filled;
        color=lightgrey;
        node [style=filled,color=white];
        a0 -> a1 -> a2 -> a3;
        label = "process #1";
      }
    
      subgraph cluster_1 {
        node [style=filled];
        b0 -> b1 -> b2 -> b3;
        label = "process #2";
        color=blue
      }
      start -> a0;
      start -> b0;
      a1 -> b3;
      b2 -> a3;
      a3 -> a0;
      a3 -> end;
      b3 -> end;
    
      start [shape=Mdiamond];
      end [shape=Msquare];
    }
    \caption{Un gráfico escrito en Dot}
    \label{fig:graficoDot}
\end{figure}

